\documentclass[11pt]{article}
\usepackage[utf8]{inputenc}
\usepackage{amsmath,amssymb,amsthm}
\usepackage[margin=2.5cm]{geometry}
\usepackage{hyperref}
\newtheorem{theorem}{Theorem}
\newtheorem{proposition}[theorem]{Proposition}
\newtheorem{conjecture}[theorem]{Conjecture}
\theoremstyle{remark}\newtheorem{remark}[theorem]{Remark}
\title{Sharp Rational Moment Hierarchies for Finite-Range Blockade Chains:\\ Kernel-Checked Results in Lean~4}
\author{FBS3R--APEX research programme}
\date{27 November 2026}
\begin{document}
\maketitle
\begin{abstract}
Let $A_{n,r}$ be the adjacency matrix of the graph whose vertices are binary words of length $n$ in which any two $1$s are more than $r$ sites apart, with edges between words differing in one site ($r=1$: the Fibonacci cube, i.e.\ the PXP model with unit couplings).
We report theorems whose proofs are checked by the Lean~4 kernel (Mathlib v4.36.0, axioms \texttt{propext}, \texttt{Classical.choice}, \texttt{Quot.sound} only).
For $r=1$ and every $m\ge1$: $\sum_n \operatorname{Tr}A_{n,1}^{2m}z^n=P_m(z)/(1-z-z^2)^{m+1}$ with $P_m\in\mathbb Z[z]$, $\deg P_m=3m-2$ exactly, sign of the leading coefficient $(-1)^{m-1}$, and $P_m\equiv(2m)!\,z^{2m-1}\pmod{1-z-z^2}$, so the denominator is sharp.
For every $r\ge0$: the second and fourth moment series are $2z/D_r^2$ and $\bigl(2zD_r+4z^2(1+\dots+z^{r-1})D_r+24z^{r+2}\bigr)/D_r^3$, $D_r=1-z-z^{r+1}$.
We also certify the eighth PXP moment for all $n$ by a proved bit-mask dynamic programme, and the dominant-root asymptotics $|C_{n,r}|\sim B_r(\rho_r)\rho_r^{-n}/K_r$.
We state precisely what remains conjectural.
\end{abstract}

\section{Setting}
$C_{n,r}$ is the set of admissible words; $a_{n,r}=|C_{n,r}|$; $D_r(z)=1-z-z^{r+1}$; $B_r(z)=1+z+\dots+z^r$.
All statements below are Lean theorems in \texttt{APEX/Lean} (registry \texttt{APEX/Audit/ACCEPTED\_THEOREMS.txt}); names are given in brackets.

\section{Dimensions and the dominant root}
\begin{theorem}[\texttt{a\_of\_le}, \texttt{a\_rec}, \texttt{gf\_coeff}]
$a_{n,r}=n+1$ for $n\le r$ and $a_{n,r}=a_{n-1,r}+a_{n-r-1,r}$; equivalently $\sum_n a_{n,r}z^n=B_r(z)/D_r(z)$.
\end{theorem}
\begin{theorem}[\texttt{exists\_unique\_root}, \texttt{complex\_root\_norm}, \texttt{Br\_rho}, \texttt{tendsto\_a\_mul\_rho\_pow}]
$D_r$ has a unique root $\rho_r\in(0,1)$; it is simple, and every other complex root has modulus $>\rho_r$.
$B_r(\rho_r)=\rho_r^{-r}$, and with $K_r=\rho_r+(r+1)\rho_r^{r+1}$ one has $a_{n,r}\rho_r^n\to B_r(\rho_r)/K_r$.
For $r=1$, $\rho_1=\varphi^{-1}$ and $2\rho_1^2/K_1=1-1/\sqrt5$.
\end{theorem}
The asymptotic statement is proved without partial fractions: $u_n=a_{n,r}\rho_r^n$ satisfies an averaging recurrence $u_{n+r+1}=pu_{n+r}+qu_n$ with $p+q=1$, which contracts windows of length $r+1$ and preserves an explicit invariant.

\section{The PXP hierarchy ($r=1$)}
\begin{theorem}[\texttt{pxp\_moment\_rational}, \texttt{pxp\_moment\_residue}, \texttt{pxp\_denominator\_sharp}, \texttt{pxp\_numerator\_natDegree}]
For every $m\ge1$ there is $P_m\in\mathbb Z[z]$ with $\bigl(\sum_n\operatorname{Tr}A_{n,1}^{2m}z^n\bigr)(1-z-z^2)^{m+1}=P_m$, $\deg P_m=3m-2$, $(-1)^{m-1}[z^{3m-2}]P_m\ge1$, and $P_m\equiv(2m)!z^{2m-1}\bmod(1-z-z^2)$.
\end{theorem}
The proof counts closed walks as arrays of time trajectories (one per site) with pairwise disjoint flip sets, and runs a strong induction on the set of used time labels, carrying three predicates (degree, residue, top coefficient).
\begin{theorem}[\texttt{MaskDP.M\_eq\_Mnat}, \texttt{eighth\_moment\_gf}]
For every period $k$, $\operatorname{Tr}A_{n,1}^k$ equals the output of an explicit bit-mask dynamic programme (proved correct). Evaluated by the kernel for $k=8$, $n\le10$, it yields
$P_4=2z+22z^2+716z^3+1146z^4+8236z^5-206z^6+13114z^7-8308z^8+3152z^9-568z^{10}$ for all $n$.
\end{theorem}

\section{Universal low moments for every radius}
\begin{theorem}[\texttt{second\_moment\_series}, \texttt{fourth\_moment}]
For every $r\ge0$:
\[\sum_n\operatorname{Tr}A_{n,r}^2z^n=\frac{2z}{D_r^2},\qquad
\sum_n\operatorname{Tr}A_{n,r}^4z^n=\frac{2zD_r+4z^2(1+\dots+z^{r-1})D_r+24z^{r+2}}{D_r^3}.\]
\end{theorem}
The fourth moment uses the decomposition $\operatorname{Tr}A^4=2e_n+12s_2(n)+2h_n$ into edges, pairs of occupied sites and non-closing ``cherries'', together with the gap identities $D_r\sum_{g<r}z^g+z^rB_r=1$ and $D_r\sum_{g=r}^{2r-1}z^g+z^{2r}B_r=z^r$ (\texttt{ident\_a}, \texttt{ident\_b}).
Note that the numerator of the fourth moment has degree $2r+2$ with leading coefficient $-4$ for every $r\ge1$.

\section{Conjectures (numerical evidence only)}
\begin{conjecture}[\texttt{ConjGeneral}]
For all $m,r\ge1$, $\sum_n\operatorname{Tr}A_{n,r}^{2m}z^n=P_{m,r}/D_r^{m+1}$ with $\deg P_{m,r}=(2r+1)(m-1)+1$, leading coefficient $(-1)^{m-1}W_m$ independent of $r$ ($W_m=2,4,36,568,14560$ for $m\le5$, $r=1$), and $P_{m,r}\equiv(2m)!\,z^{(r+1)m-r}\pmod{D_r}$.
\end{conjecture}
Exact integer checks: $r\le4$, $m\le3$, $n\le22$ (\texttt{blockade\_verify.py}, polynomial division over $\mathbb Z$). The conjecture would imply the Gaussian limit of the empirical spectral measures with variance $\sigma_r^2=2\rho_r^{r+1}/K_r$, which is not proved here.

\section{Status and novelty}
Rationality of such series by transfer matrices is standard. The exact degree, the residue, the $r$-uniform fourth moment and the kernel certificates are, to our knowledge from limited searches, not in the literature; we make no priority claim. The proofs are machine-checked; the literature comparison is not complete.

\paragraph{Reproduction.} \texttt{cd APEX/Lean \&\& lake build \&\& bash ../ci\_check.sh}.
\end{document}
